\section{Related work}\label{sec:literature}

The paper sits at the intersection of three literatures: macro-fiscal
analyses of AI, tax-benefit microsimulation of technology shocks, and the
factor-share literature that motivates the experiment's central mechanism.

\textbf{AI and the public finances.} \citet{iselinnunn2026} is the direct
antecedent: they translate the \citet{karger2026} forecaster survey into
2030 growth, factor-share and wage-distribution scenarios and score them
through a federal individual income and payroll tax microsimulation built on
the IRS Public Use File, adding a corporate revenue wedge outside the
microsimulation. Their central finding --- that growth accruing to capital
raises less revenue than the same growth at fixed factor shares, by roughly a
factor of two in the Rapid scenario --- is the claim this paper re-derives on
independent machinery and extends to transfers, poverty and states.
\citet{karger2026} elicit forecasts from economists, AI-company staff, policy
researchers, superforecasters and the public; the median expects modest
GDP acceleration by 2030 with substantial disagreement about wage
distribution effects --- disagreement this paper propagates rather than
resolves. \citet{acemoglu2025} argues on task-level grounds for modest
aggregate effects; \citet{bastani2024} survey the tax-design questions AI
raises, including whether existing bases can reach AI-driven rents.

\textbf{Microsimulation of AI shocks.} Closest in method are two companion
studies on European systems. \citet{doorley2026} run scenario-based AI
employment and wage shocks through the Irish SWITCH model, and
\citet{ahmadi2026} runs occupation-level exposure-based displacement, wage and
capital shocks through PolicyEngine UK, treating the incidence of
displacement as an explicit scenario axis. This paper is the United States
member of that family, but it differs in the shock design: rather than
displacement incidence, it varies the factor composition and
wage-distribution shape of income growth, matching the Budget Lab's
calibration to enable the cell-level reconciliation of
Section~\ref{sec:reconciliation}. Occupation-level exposure measures
\citep{felten2021, eloundou2024} are the natural next refinement of the wage
channel here, as Section~\ref{sec:discussion} discusses; both this paper's
repository and the Budget Lab's roadmap point the same direction.

\textbf{Factor shares and automation.} The premise that technology can move
income between labor and capital at scale is not hypothetical: the global
decline of the labor share is documented by \citet{karabarbounisneiman2014},
and \citet{acemoglurestrepo2018} provide the task-model mechanics by which
automation reallocates value added from labor to capital. What
microsimulation adds to that literature is the statutory detail: a dollar
moved from wages to capital income changes its withholding, its marginal
rate schedule, its payroll treatment, its state treatment, and the transfer
payments of the household it leaves --- none of which aggregate models
carry. The scenarios here take the CBO baseline as the no-AI counterfactual,
using the same February 2026 projections the Budget Lab anchors to
\citep{cbo2026}.
